% ==========================================================================
% Section 8: Conclusion and Future Work
% ==========================================================================
\section{Conclusion and Future Work}
\label{sec:conclusion-future-work}

\subsection{Summary of the Implemented System}
\label{sec:summary-implemented-system}

L\_RAG combines a Vietnamese legal chunk corpus, dense FAISS retrieval,
optional BM25/RRF fusion, graph expansion, cross-encoder reranking, grounded
generation, and saved per-question evaluation records. The ablation report
isolates five reported component analyses---embedding text, retrieval mode,
graph policy, reranking, and generator prompting---grouped into four run
families and evaluated with fixed question sets and paired metrics.

\subsection{Main Contributions}
\label{sec:main-contributions}

The empirical contribution is a traceable set of controlled decisions: the
embedded text representation changes document success but not exact chunk
coverage; the tested hybrid and graph variants impose measurable coverage
losses; reranking within a common pool produces substantial ranking gains; and
CoT does not improve lexical answer overlap under fixed context. The report
also separates retrieval evidence, answer overlap, refusal formatting,
citation structure, and efficiency instead of collapsing them into one score.

\subsection{Lessons Learned}
\label{sec:lessons-learned}

The central lesson is that component labels are not outcomes. Metadata,
hybrid retrieval, graph expansion, CoT, and cited output are not quality
guarantees. Candidate control and paired question-level deltas reveal effects
obscured by aggregate scores: metadata can improve document discovery while
delaying exact chunks, graph expansion increases density while removing
evidence, and rerankers can only improve what has entered C30. The graph loss
also reflects a replacement policy that starts from five seeds and retains at
most ten expanded chunks; preserving the original seeds is the next targeted
test. Frozen manifests and complete row alignment are essential before
interpreting a difference as causal.

\subsection{Recommended System Configuration}
\label{sec:recommended-configuration}

For the evaluated benchmark, the recommended evidence path is dense retrieval
with the Jina reranker when its latency is acceptable, a 10-chunk final
context, and the Base prompt. Jina is promising within the tested dense--BM25
RRF $C_{30}$ pool; its performance on a dense-only candidate pool remains
unmeasured. Use the graph branch only after seed-preserving expansion is
validated, and do not select the current hybrid branch as a default. Its
reported latency includes amortized sharded-BM25 batch computation, so online
latency must be measured separately. This is a benchmark-bounded configuration,
not a deployment approval. A real deployment requires corpus-validity checks,
source display, monitored latency, privacy controls, abstention handling, and
legal expert review.

\subsection{Future Improvements}
\label{sec:future-improvements}

Priority work is (1) blinded legal-expert annotation for correctness,
completeness, and claim-to-citation entailment; (2) seed-preserving graph and
neighbor-budget sweeps; (3) BM25/tokenization and RRF-weight tuning under
full-online latency measurement; (4) candidate-pool sweeps beyond C30; and
(5) temporal-validity and external-reference evaluation. These studies should
retain per-case candidates and provenance so failures can be audited.

\subsection{Final Conclusion}
\label{sec:final-conclusion}

The ablation evidence supports a practical conclusion: improve the ordering of
available legal evidence with reranking before adding unvalidated retrieval
complexity. Jina is the strongest observed reranker in the fixed-pool study,
while the current hybrid, graph, and CoT variants do not improve their primary
quality measures enough to offset their costs or risks. The result is an
evidence-based engineering recommendation, bounded by the benchmark and not a
claim of legal correctness.
